\documentclass[11pt,twocolumn]{article}
\usepackage[a4paper,margin=2cm]{geometry}
\usepackage{fontspec}
\begin{document}
\twocolumn[{\section*{Spaced Practice and Delayed Recall}}]
\section*{Introduction}
Learners often reread a chapter several times in a single evening and feel confident the next morning. That confidence fades quickly. Spaced practice distributes the same effort across days, so that each review happens just as the memory begins to weaken. This short report describes a small classroom study of that effect and the conditions under which it held. We focus on delayed recall, because it is the outcome that matters when the material is needed again weeks later.

\section*{Method}
Forty learners studied twelve concepts from an introductory statistics course. Half of them reviewed each concept three times on the first day. The other half reviewed each concept once on the first day, once two days later and once a week later. Every review session lasted about twenty minutes, and the total study time was identical in both groups. Nobody was told which schedule was expected to work better, and the same instructor taught both groups.

\section*{Results}
One week after the last session, the spaced group recalled 78 percent of the concepts and the massed group recalled 52 percent. After a month the gap was even larger. The difference was consistent across the easier and the harder concepts, although the harder concepts benefited slightly more. Learners in the spaced group also rated their confidence more accurately, which suggests that the schedule improved their judgement as well as their memory.

\section*{Discussion}
The finding is not new, but it is easy to forget when planning a course. A practical study tool should therefore schedule reviews automatically, keep the effort per session small and show learners how their recall develops over time. Future work will test the same design with longer texts, with mixed Persian and English material and with learners who study on their own.

\section*{Limitations}
The sample was small and came from a single course, so the size of the effect should not be generalised without care. The concepts were short definitions rather than long arguments, and recall was measured with the same kind of questions that were used during practice. A stronger design would also measure transfer to new problems, and it would follow learners for a full semester rather than a single month.

\section*{Conclusion}
Distributing the same amount of practice over several days produced clearly better delayed recall than concentrating it in one evening. The schedule cost nothing extra in study time. For a learning environment that aims at durable knowledge, spaced review should be the default rather than an option that learners must discover for themselves.

\end{document}
